\documentclass[aspectratio=169]{beamer}
\usetheme{metropolis}
\usepackage{appendixnumberbeamer}
\usepackage{booktabs, hyperref}

\input{mgmt675-style}

\usepackage{tikz}
\usetikzlibrary{shapes.geometric, arrows.meta, positioning, calc}

% Title info
\subtitle{MGMT 675: Generative AI for Finance}
\title{Module 7: Verification, Auditing, and Governance}
\author{Kerry Back}
\date{}

\begin{document}

\maketitle

% ========================================
% SECTION 1: THE TRUST SPECTRUM
% ========================================
\section{The Trust Spectrum}

\begin{frame}[shrink=5]{AI with Code Execution: A Fast Junior Analyst}
AI with code execution is like a \alert{fast junior analyst} --- capable but needs oversight. The level of oversight depends on the stakes.
\vspace{0.3cm}
\begin{columns}[T]
\begin{column}{0.33\textwidth}
\begin{shadedbox}[title=\textbf{Low-Stakes}]
\begin{itemize}\small
  \item Exploratory charts
  \item Directional trends
  \item Internal brainstorming
\end{itemize}
\vspace{0.2cm}
\centering\small\textit{Quick sanity check.}
\end{shadedbox}
\end{column}
\begin{column}{0.33\textwidth}
\begin{shadedbox}[title=\textbf{Medium-Stakes}]
\begin{itemize}\small
  \item Team tools and skills
  \item Recurring reports
  \item Internal dashboards
\end{itemize}
\vspace{0.2cm}
\centering\small\textit{Test with known answers.}
\end{shadedbox}
\end{column}
\begin{column}{0.33\textwidth}
\begin{shadedbox}[title=\textbf{High-Stakes}]
\begin{itemize}\small
  \item Client deliverables
  \item Regulatory filings
  \item Board presentations
\end{itemize}
\vspace{0.2cm}
\centering\small\textit{Full verification protocol.}
\end{shadedbox}
\end{column}
\end{columns}
\vspace{0.3cm}
The more consequential the decision, the more you verify. \alert{Not distrust --- calibrated confidence.}
\end{frame}

% ========================================
% SECTION 2: SILENT ANALYTICAL ERRORS
% ========================================
\section{AI Makes Silent Analytical Errors}

\begin{frame}{The Real Risk: Not Hallucination --- Bugs}
AI in code-execution mode doesn't invent facts the way a chatbot might. Instead, it makes \alert{silent analytical errors}: wrong filters, dropped rows, incorrect joins.
\vspace{0.5cm}
\begin{columns}[T]
\begin{column}{0.5\textwidth}
\textbf{Common errors}
\begin{itemize}\small
  \item Missing data silently dropped
  \item Date or currency parsing surprises
  \item Filters slightly off; aggregating before vs.\ after filtering
\end{itemize}
\end{column}
\begin{column}{0.5\textwidth}
\textbf{Why it's dangerous}
\begin{itemize}\small
  \item The output \textit{looks} professional
  \item Charts render cleanly and nobody questions a polished result
  \item \alert{Confidence without verification is the real risk}
\end{itemize}
\end{column}
\end{columns}
\end{frame}

\begin{frame}{Catching a Silent Error}
\begin{columns}[T]
\begin{column}{0.5\textwidth}
\textbf{The prompt}\\[0.3em]
\small\textit{``What's the average profit margin by sub-category? Exclude any sub-categories with negative total profit.''}
\vspace{0.3cm}

\textbf{The trap}
\begin{itemize}\small
  \item Does AI filter \textit{before} or \textit{after} aggregating?
  \item Does it drop rows with negative profit, or sub-categories with negative \textit{total} profit?
  \item The order matters --- and AI often gets it wrong
\end{itemize}
\end{column}
\begin{column}{0.5\textwidth}
\textbf{How to catch it}
\begin{itemize}\small
  \item \alert{``Show me the code''} --- inspect the filter logic
  \item \alert{``How many rows are in the result?''} --- does it match your expectation?
  \item \alert{``List the excluded sub-categories''} --- are they the right ones?
\end{itemize}
\vspace{0.3cm}
\textbf{Lesson}: The more specific the filter logic, the more likely AI misinterprets it. Always verify multi-step operations.
\end{column}
\end{columns}
\end{frame}

% ========================================
% SECTION 3: THE VERIFICATION CHECKLIST
% ========================================
\section{The Verification Checklist}

\begin{frame}{A 4-Step Verification Protocol}
Apply this to \alert{any} AI-generated analysis.
\vspace{0.5cm}
\begin{enumerate}
  \item \textbf{Sanity-check with known answers}\\
  \small Ask questions where you already know the result: row count, column names, a total you can verify in Excel.
  \item \textbf{Ask AI to show its methodology}\\
  \small ``Show me the code'' or ``Show me the SQL.'' Read the logic, not just the output.
  \item \textbf{Cross-check by rephrasing}\\
  \small Ask the same question differently. Start a fresh conversation and ask again. Do answers agree?
  \item \textbf{Spot-check edge cases}\\
  \small Test the smallest group, zero values, null fields, boundary dates. Errors hide at the edges.
\end{enumerate}
\end{frame}

\begin{frame}{}
\vspace{2cm}
\begin{center}
{\Large\bfseries Trust is a process, not a setting.}
\vspace{1cm}

Every AI output deserves the same scrutiny\\you'd give a new analyst's first deliverable.
\end{center}
\end{frame}

% ========================================
% SECTION 4: VALIDATING WORK FROM MODULES 1--4
% ========================================
\section{Hands-On: Validating Prior Work}

\begin{frame}{Reproducibility Test}
Take analysis from earlier modules and test whether it reproduces.
\vspace{0.3cm}
\begin{columns}[T]
\begin{column}{0.5\textwidth}
\textbf{The exercise}
\begin{enumerate}\small
  \item Pick a financial analysis from Module 2 (portfolio) or Module 4 (DCF)
  \item Start a \alert{fresh conversation}
  \item Give the same data and ask the same question
  \item Compare: do the results match?
\end{enumerate}
\end{column}
\begin{column}{0.5\textwidth}
\textbf{What you're testing}
\begin{itemize}\small
  \item \textbf{Reproducibility}: Same data, same question --- same answer?
  \item \textbf{Methodology}: Did AI use the same approach both times?
  \item \textbf{Stability}: Small rephrasing shouldn't change the result
\end{itemize}
\end{column}
\end{columns}
\vspace{0.3cm}
If the numbers don't match, that's not a failure --- it's \alert{exactly why we verify}. And it's exactly why \textbf{skills} (Module 3) matter.
\end{frame}

\begin{frame}{Auditing a DCF Model}
Take the DCF built in Module 4. Start a fresh conversation. Ask Claude to \alert{critique} it.
\vspace{0.3cm}
\begin{columns}[T]
\begin{column}{0.5\textwidth}
\textbf{What to ask}
\begin{itemize}\small
  \item ``Identify 3+ questionable assumptions''
  \item ``Check formula consistency --- does FCF match the pro formas?''
  \item ``Is the terminal value reasonable? Compare implied terminal growth to GDP growth.''
\end{itemize}
\end{column}
\begin{column}{0.5\textwidth}
\textbf{Why this matters}
\begin{itemize}\small
  \item The most valuable use of AI is often having it \alert{check} work, not produce it
  \item Cross-evaluation: use a second conversation to critique the first
  \item Compare original vs.\ revised valuations
\end{itemize}
\end{column}
\end{columns}
\end{frame}

\begin{frame}{Red-Teaming a Deployed Skill}
A database skill is the kind of tool people \alert{stop questioning} once deployed. Test it like any tool your team depends on.
\vspace{0.3cm}
\begin{columns}[T]
\begin{column}{0.5\textwidth}
\textbf{Adversarial queries}
\begin{enumerate}\small
  \item Query for nonexistent data (``Sales in 2030'')
  \item Request something not in the schema (``Customer satisfaction scores'')
  \item Use ambiguous language (``Best employee'')
\end{enumerate}
\end{column}
\begin{column}{0.5\textwidth}
\textbf{For each failure}
\begin{itemize}\small
  \item \textbf{Diagnose}: Bad SQL? Missing context? Wrong assumption?
  \item \textbf{Fix}: Add clarification, examples, or constraints to SKILL.md
  \item \textbf{Retest}: Run the query again after the fix
\end{itemize}
\vspace{0.3cm}
This is how production skills get hardened. \alert{Every failure makes the skill better.}
\end{column}
\end{columns}
\end{frame}

% ========================================
% SECTION 5: RED-TEAMING PROTOCOLS
% ========================================
\section{Red-Teaming Protocols}

\begin{frame}{Systematic Red-Teaming}
Red-teaming goes beyond ad-hoc testing.  A \alert{protocol} ensures you test every failure mode systematically.
\vspace{0.3cm}
\begin{center}\small
\begin{tabular}{@{} l l @{}}
\toprule
\textbf{Test Category} & \textbf{What You Try} \\
\midrule
Out-of-scope questions & Ask about data that doesn't exist in the schema \\
Conflicting documents & Upload docs with contradictory numbers \\
Edge cases & Zero values, null fields, single-row groups, boundary dates \\
Hand-written SQL comparison & Write the query yourself and compare to AI's result \\
Ambiguous language & ``Best customer'' --- by revenue? frequency? margin? \\
\bottomrule
\end{tabular}
\end{center}
\vspace{0.3cm}
\alert{A skill that hasn't been red-teamed is a liability, not an asset.}
\end{frame}

\begin{frame}[shrink=5]{Red-Teaming a DCF Model}
Apply red-teaming protocols to the DCF models built in Module 4.
\vspace{0.2cm}
\begin{columns}[T]
\begin{column}{0.5\textwidth}
\begin{shadedbox}[title=\textbf{Financial Red-Team Tests}]
\begin{enumerate}\small
  \item Set revenue growth to 0\% --- does the model break?
  \item Use a WACC below the risk-free rate --- does AI flag it?
  \item Feed a company with negative FCF and compare AI's output to a hand-built Excel model
\end{enumerate}
\end{shadedbox}
\end{column}
\begin{column}{0.5\textwidth}
\begin{shadedbox}[title=\textbf{What You Learn}]
\begin{itemize}\small
  \item Where the model silently produces nonsense
  \item Which assumptions the AI accepts without questioning
  \item How to improve guardrails so the verification checklist catches the errors
\end{itemize}
\end{shadedbox}
\end{column}
\end{columns}
\end{frame}

% ========================================
% SECTION 6: AUDIT LOGGING
% ========================================
\section{Audit Logging}

\begin{frame}{JSONL Audit Logs}
Every AI interaction in a production system should produce an \alert{audit log entry}. The standard format is JSONL (one JSON object per line).
\vspace{0.3cm}
\begin{columns}[T]
\begin{column}{0.5\textwidth}
\begin{shadedbox}[title=\textbf{What to Log}]
\begin{itemize}\small
  \item \textbf{Identity}: timestamp, user, model
  \item \textbf{Content}: prompt and response
  \item \textbf{Metadata}: token count, latency, cost
\end{itemize}
\end{shadedbox}
\end{column}
\begin{column}{0.5\textwidth}
\begin{shadedbox}[title=\textbf{Why It Matters}]
\begin{itemize}\small
  \item \textbf{Compliance}: Regulators can audit what the AI said
  \item \textbf{Debugging}: Trace errors back to the exact prompt
  \item \textbf{Cost and quality}: Track token usage and identify common failures for skill updates
\end{itemize}
\end{shadedbox}
\end{column}
\end{columns}
\vspace{0.3cm}
Store logs locally or in a cloud database.  \alert{If it's not logged, it didn't happen.}
\end{frame}

% ========================================
% SECTION 7: THREE LAYERS OF AUDITABILITY
% ========================================
\section{Three Layers of Auditability}

\begin{frame}[shrink=5]{Defense in Depth}
\begin{center}
\begin{tikzpicture}[
  every node/.style={font=\small},
  layer/.style={draw=titlegray, rounded corners, minimum height=0.8cm, fill=#1, text=titlegray, font=\small\bfseries, align=center}
]
\node[layer=accentblue!15, minimum width=10cm] (l1) at (0,2) {Layer 1: RAG Citations};
\node[layer=excelinput, minimum width=8cm] (l2) at (0,1) {Layer 2: Code Transparency};
\node[layer=alertorange!12, minimum width=6cm] (l3) at (0,0) {Layer 3: Audit Logs};

\node[right=0.3cm of l1, font=\scriptsize, text=titlegray, text width=3cm] {Answers grounded in source documents};
\node[right=0.3cm of l2, font=\scriptsize, text=titlegray, text width=3cm] {SQL/Python visible for review};
\node[right=0.3cm of l3, font=\scriptsize, text=titlegray, text width=3cm] {Complete trail in JSONL};
\end{tikzpicture}
\end{center}
\vspace{0.2cm}
\begin{enumerate}
  \item \textbf{RAG citations}: Ground answers in source documents with traceable references (Module 8). AI answers based on \alert{your content}, not its training data.
  \item \textbf{Code transparency}: AI displays the SQL or Python it used. Analysts review the logic, not just the answer.
  \item \textbf{Audit logs}: Every prompt, response, and token count recorded in JSONL. The system creates a \textit{complete trail} of every AI decision.
\end{enumerate}
\vspace{0.3cm}
The more layers you add, the harder it is for errors to reach your team. \alert{Same principle as cybersecurity.}
\end{frame}

% ========================================
% SECTION 6: SECURITY AND GOVERNANCE
% ========================================
\section{Security and Governance}

\begin{frame}{Data Policies: Training and Retention}
Before sending data to any AI provider, know the answers to two questions.
\vspace{0.3cm}
\begin{columns}[T]
\begin{column}{0.5\textwidth}
\textbf{Will they train on your data?}
\begin{itemize}\small
  \item \alert{Consumer tiers} (free ChatGPT, free Claude) may use your input to improve models
  \item \alert{API and enterprise tiers} typically do not
  \item Check your agreement and settings
\end{itemize}
\end{column}
\begin{column}{0.5\textwidth}
\textbf{How long do they keep it?}
\begin{itemize}\small
  \item Providers retain prompts for a window (often 30 days) for safety monitoring
  \item Enterprise contracts may offer zero retention
  \item Know the policy \alert{before} you send sensitive data
\end{itemize}
\end{column}
\end{columns}
\end{frame}

\begin{frame}[shrink=5]{The Schema-Only Pattern}
AI can analyze sensitive data \alert{without ever reading it}. It only needs the table structure.
\vspace{0.2cm}
\begin{center}
\begin{tikzpicture}[
  node distance=0.5cm and 1.2cm,
  every node/.style={font=\small},
  sbox/.style={draw=titlegray, rounded corners, minimum width=2cm, minimum height=0.7cm, fill=excelinput, text=titlegray, font=\scriptsize\bfseries, align=center},
  sarrow/.style={-{Stealth[length=2.5mm]}, thick, draw=accentblue}
]
\node[sbox] (schema) {Schema\\{\tiny (table names, columns)}};
\node[sbox, right=1.2cm of schema, fill=accentblue!15] (ai) {AI};
\node[sbox, right=1.2cm of ai] (sql) {SQL\\Code};
\node[sbox, below=0.8cm of sql, fill=alertorange!12] (local) {Runs\\Locally};
\node[sbox, below=0.8cm of schema] (db) {Your\\Database};

\draw[sarrow] (schema) -- (ai) node[midway, above, font=\scriptsize\itshape, text=titlegray] {safe to share};
\draw[sarrow] (ai) -- (sql);
\draw[sarrow] (sql) -- (local);
\draw[sarrow] (db) -- (local) node[midway, below, font=\scriptsize\itshape, text=titlegray] {data stays here};

\draw[dashed, thick, alertorange] (-1.5,-0.4) -- (2.2,-0.4) node[midway, below=0.1cm, font=\scriptsize, text=alertorange] {};
\node[font=\scriptsize, text=alertorange] at (0.3, -0.6) {on-premise boundary};
\end{tikzpicture}
\end{center}
\vspace{0.1cm}
\begin{columns}[T]
\begin{column}{0.5\textwidth}
\textbf{How it works}
\begin{enumerate}\small
  \item Describe tables in instructions: names, columns, types
  \item Ask AI to write the query
  \item AI writes SQL from the schema alone; the script runs locally and output stays local
\end{enumerate}
\end{column}
\begin{column}{0.5\textwidth}
\textbf{Why it's compliant}
\begin{itemize}\small
  \item Table metadata is \alert{not PII} --- safe to share
  \item Sensitive records stay on your server
  \item AI produces the \textit{tool}, not the \textit{output}
  \item Works for FERPA, HIPAA, SOX, GDPR
\end{itemize}
\end{column}
\end{columns}
\end{frame}

% ========================================
% SECTION 7: ADOPTION ROADMAP
% ========================================
\section{AI Adoption Roadmap}

\begin{frame}{The Sandbox-Audit-Deploy Pattern}
Before deploying any AI workflow to your team, follow three phases.
\vspace{0.3cm}
\begin{center}
\begin{tikzpicture}[
  node distance=1.5cm,
  every node/.style={font=\small},
  sbox/.style={draw=titlegray, rounded corners, minimum width=3.2cm, minimum height=1cm, fill=excelinput, text=titlegray, font=\small\bfseries, align=center},
  sarrow/.style={-{Stealth[length=3mm]}, thick, draw=accentblue}
]
\node[sbox] (sandbox) {1. Sandbox\\{\scriptsize\normalfont Test with dummy data}};
\node[sbox, right=of sandbox] (audit) {2. Audit\\{\scriptsize\normalfont Review logs and logic}};
\node[sbox, right=of audit] (deploy) {3. Deploy\\{\scriptsize\normalfont Roll out with monitoring}};

\draw[sarrow] (sandbox) -- (audit);
\draw[sarrow] (audit) -- (deploy);
\end{tikzpicture}
\end{center}
\vspace{0.3cm}
Test with synthetic data, review logs and logic, then roll out with weekly monitoring.
\end{frame}

\begin{frame}{Adoption Timeline: From Quick Wins to Strategic}
Each adoption stage maps to a trust level.
\vspace{0.3cm}
\begin{enumerate}
  \item \textbf{Quick wins} (weeks) --- \textit{light verification}:
  \begin{itemize}\small
    \item Chat for research, summarization, drafts
    \item Upload spreadsheets for ad-hoc analysis
  \end{itemize}
  \item \textbf{Medium-term} (months) --- \textit{test-before-deploy}:
  \begin{itemize}\small
    \item Skills for recurring workflows (Module 3)
    \item MCP integrations and Streamlit apps (Modules 7, 8)
  \end{itemize}
  \item \textbf{Strategic} (quarters) --- \textit{full audit trail}:
  \begin{itemize}\small
    \item Production apps with authentication
    \item AI usage policy and governance framework
  \end{itemize}
\end{enumerate}
\end{frame}

% ========================================
% SECTION 8: SMALL LANGUAGE MODELS
% ========================================
\section{Small Language Models and Fine-Tuning}

\begin{frame}{Improving Reliability: Fine-Tuning and SLMs}
So far we've verified AI output \textit{after the fact}.  A complementary approach: make the model itself more reliable for your specific tasks.
\vspace{0.3cm}
\begin{columns}[T]
\begin{column}{0.5\textwidth}
\textbf{Fine-tuning an LLM}
\begin{itemize}\small
  \item Train a general-purpose model on your domain data (10-K filings, internal memos, deal documents)
  \item The model learns your terminology, formats, and conventions
  \item Fewer silent errors on domain-specific tasks
\end{itemize}
\end{column}
\begin{column}{0.5\textwidth}
\textbf{Small Language Models (SLMs)}
\begin{itemize}\small
  \item Models with $\sim$1--8 billion parameters (vs.\ hundreds of billions for LLMs)
  \item Can be fine-tuned in hours on a single GPU
  \item Run on-premise or on-device --- \alert{data never leaves your network}
\end{itemize}
\end{column}
\end{columns}
\vspace{0.3cm}
\alert{A specialized small model can outperform a general-purpose large model on narrow, repetitive tasks.}
\end{frame}

\begin{frame}{Why SLMs for Agentic Systems?}
Belcak et al.\ (2025), ``\href{https://arxiv.org/abs/2506.02153}{Small Language Models are the Future of Agentic AI}'':
\vspace{0.3cm}
\begin{itemize}
  \item Most agent subtasks are \alert{repetitive, scoped, and non-conversational} --- classifying documents, extracting fields, formatting outputs
  \item SLMs fine-tuned on a single task outperform LLMs on format compliance and consistency
  \item 10--30$\times$ cheaper in latency, energy, and compute vs.\ large models
  \item On-device deployment enables real-time, offline inference with full data privacy
  \item Agents can use \alert{multiple models}: an LLM for complex reasoning, SLMs for routine subtasks
\end{itemize}
\end{frame}

\begin{frame}{From LLM to SLM: The Conversion Process}
The paper proposes a six-step migration path:
\vspace{0.3cm}
\begin{enumerate}
  \item \textbf{Log} agent interactions (prompts, responses, tool calls)
  \item \textbf{Curate} training data --- remove PII, filter for quality
  \item \textbf{Cluster} requests to identify recurring task patterns
  \item \textbf{Select} an SLM sized to each task's complexity
  \item \textbf{Fine-tune} using parameter-efficient methods (LoRA) --- typically 10K--100K examples suffice
  \item \textbf{Repeat} --- retrain as usage patterns evolve
\end{enumerate}
\vspace{0.3cm}
\alert{The audit logs from earlier in this module are exactly the data you need for step 1.}
\end{frame}

% ========================================
% SECTION 9: SUMMARY
% ========================================
\section{Summary}

\begin{frame}{Summary}
\begin{columns}[T]
\begin{column}{0.25\textwidth}
\begin{shadedbox}[title=\textbf{Verification}]
\begin{itemize}\scriptsize
  \item Trust spectrum
  \item 4-step protocol
  \item Silent errors
\end{itemize}
\end{shadedbox}
\end{column}
\begin{column}{0.25\textwidth}
\begin{shadedbox}[title=\textbf{Red-Teaming}]
\begin{itemize}\scriptsize
  \item Systematic protocols
  \item Edge cases
  \item SQL comparison
\end{itemize}
\end{shadedbox}
\end{column}
\begin{column}{0.25\textwidth}
\begin{shadedbox}[title=\textbf{Audit Logs}]
\begin{itemize}\scriptsize
  \item JSONL format
  \item Compliance trail
  \item Cost tracking
\end{itemize}
\end{shadedbox}
\end{column}
\begin{column}{0.25\textwidth}
\begin{shadedbox}[title=\textbf{Governance}]
\begin{itemize}\scriptsize
  \item Data policies
  \item Schema-only
  \item Adoption roadmap
\end{itemize}
\end{shadedbox}
\end{column}
\end{columns}
\vspace{0.5cm}
\begin{center}
\alert{Verify, govern, then deploy.}
\end{center}
\end{frame}

\end{document}
